% !TEX program = xelatex
\documentclass[12pt,a4paper]{ctexart}

% ============ Packages ============
\usepackage{amsmath,amssymb,mathtools,bm,cancel}
\usepackage{amsthm}
\usepackage[colorlinks=true,linkcolor=blue,citecolor=blue]{hyperref}
\usepackage{booktabs}
\usepackage{enumitem}
\usepackage[margin=2.5cm]{geometry}
\usepackage{xcolor}
\usepackage{fancyhdr}
\usepackage{titlesec}
\usepackage{listings}
\usepackage{graphicx}

% ============ Theorems ============
\theoremstyle{definition}
\newtheorem{definition}{定义}[section]
\newtheorem{theorem}{定理}[section]
\newtheorem{lemma}[theorem]{引理}
\newtheorem{corollary}[theorem]{推论}
\newtheorem{proposition}[theorem]{命题}
\newtheorem{remark}{注记}[section]
\theoremstyle{plain}
\newtheorem*{proof*}{证明}

% ============ Math Notation ============
\newcommand{\R}{\mathbb{R}}
\newcommand{\E}{\mathbb{E}}
\newcommand{\diag}{\operatorname{diag}}
\newcommand{\tr}{\operatorname{tr}}
\newcommand{\vect}{\operatorname{vec}}
\newcommand{\eps}{\varepsilon}
\newcommand{\grad}{\nabla}
\newcommand{\norm}[1]{\left\lVert #1 \right\rVert}
\newcommand{\ip}[2]{\left\langle #1,#2 \right\rangle}
\newcommand{\Id}{I}
\DeclareMathOperator*{\argmin}{arg\,min}
\DeclareMathOperator*{\argmax}{arg\,max}

% ============ Code ============
\lstset{
  basicstyle=\ttfamily\small,
  backgroundcolor=\color{gray!5},
  frame=single,
  framerule=0.5pt,
  rulecolor=\color{gray!40},
  breaklines=true,
  numbers=left,
  numbersep=6pt,
  numberstyle=\tiny\color{gray},
  tabsize=2,
  keywordstyle=\color{blue!70!black},
  commentstyle=\color{gray!50!black},
  stringstyle=\color{orange!70!black},
  showstringspaces=false,
}

% ============ Formatting ============
\pagestyle{fancy}
\fancyhf{}
\fancyhead[L]{第一周：GD/SGD 理论基石 —— 习题解答}
\fancyhead[R]{\thepage}
\fancyfoot[C]{上海财经大学·计算机机械社}
\renewcommand{\headrulewidth}{0.5pt}

\titleformat{\section}{\Large\bfseries}{}{0em}{}
\titleformat{\subsection}{\large\bfseries}{}{0em}{}

% ============ Document ============
\begin{document}

\begin{center}
\Large\textbf{第一周：GD 与 SGD 理论基石}\\[4pt]
\textbf{—— 理论与工程练习完整解答 ——}

\vspace{0.5cm}
\normalsize 上海财经大学·计算机机械社
\end{center}

\vspace{0.3cm}
\rule{\textwidth}{0.5pt}

% ======================================================================
\section{理论练习}
% ======================================================================

% ---------- 练习 1 ----------
\subsection{练习 1：下降引理的逐行推导}

\textbf{题目：} 若 $\norm{\grad f(x) - \grad f(y)} \le L\norm{x-y}$，证明
$f(y) \le f(x) + \ip{\grad f(x)}{y-x} + \frac{L}{2}\norm{y-x}^2$。

\begin{proof*}[解答]
设 $x, y \in \R^d$。考虑线段 $x + t(y-x)$（$t \in [0,1]$），定义辅助函数
$h(t) = f(x + t(y-x))$。则 $h(0) = f(x)$，$h(1) = f(y)$。

由链式法则：
\begin{equation}
  h'(t) = \big\langle \grad f(x + t(y-x)),\; y - x \big\rangle.
\end{equation}

应用微积分基本定理：
\begin{equation}\label{s1:step1}
  f(y) - f(x) = h(1) - h(0)
  = \int_{0}^{1} h'(t)\,dt
  = \int_{0}^{1} \big\langle \grad f(x + t(y-x)),\; y - x \big\rangle\,dt.
\end{equation}

将内积拆分为两项（加减 $\grad f(x)$）：
\begin{equation}\label{s1:step2}
  \begin{aligned}
    f(y) - f(x)
    &= \int_{0}^{1} \ip{\grad f(x)}{y-x}\,dt
       + \int_{0}^{1} \ip{\grad f(x + t(y-x)) - \grad f(x)}{y-x}\,dt \\[4pt]
    &= \ip{\grad f(x)}{y-x}
       + \int_{0}^{1} \ip{\grad f(x + t(y-x)) - \grad f(x)}{y-x}\,dt.
  \end{aligned}
\end{equation}

第一项不含 $t$，故积分为简单的 $\ip{\grad f(x)}{y-x}$。

对第二项应用 Cauchy-Schwarz 不等式：
\begin{equation}\label{s1:step3}
  \big|\ip{\grad f(x + t(y-x)) - \grad f(x)}{y-x}\big|
  \le \norm{\grad f(x + t(y-x)) - \grad f(x)} \cdot \norm{y-x}.
\end{equation}

利用 $L$-光滑性假设：$\norm{\grad f(x + t(y-x)) - \grad f(x)} \le L \norm{t(y-x)} = L t \norm{y-x}$。代入：
\begin{equation}\label{s1:step4}
  \big|\ip{\grad f(x + t(y-x)) - \grad f(x)}{y-x}\big|
  \le L t \norm{y-x}^2.
\end{equation}

将上式代入式~(\ref{s1:step2}) 的积分：
\begin{equation}\label{s1:step5}
  \begin{aligned}
    f(y) - f(x)
    &\le \ip{\grad f(x)}{y-x}
       + \int_{0}^{1} L t \norm{y-x}^2\,dt \\[4pt]
    &= \ip{\grad f(x)}{y-x}
       + L \norm{y-x}^2 \int_{0}^{1} t\,dt \\[4pt]
    &= \ip{\grad f(x)}{y-x} + \frac{L}{2}\norm{y-x}^2.
  \end{aligned}
\end{equation}

移项即得到下降引理：
\begin{equation}
  \boxed{f(y) \le f(x) + \ip{\grad f(x)}{y-x} + \frac{L}{2}\norm{y-x}^2.}
\end{equation}

\textbf{注记：} 该不等式将函数值 $f(y)$ 的上界表示为一个线性项（一阶近似）
加上一个二次罚项（由光滑性常数 $L$ 度量）。正是这个二次界保证了：
只要步长 $\eta < 2/L$，梯度下降每步一定会使函数值下降。

\end{proof*}

% ---------- 练习 2 ----------
\subsection{练习 2：GD 在二次函数上的精确收敛因子}

\textbf{题目：} 考虑 $f(\theta) = \frac12\theta^\top A\theta$，$A = \diag(\lambda_1, \ldots, \lambda_d)$。
写出 GD 在第 $i$ 个坐标上的递推，并推导最优步长 $\eta$ 和收敛因子 $(\kappa - 1)/(\kappa + 1)$。

\begin{proof*}[解答]
\noindent\textbf{第一步：按坐标解耦。}

梯度为 $\grad f(\theta) = A\theta = (\lambda_1\theta_1, \dots, \lambda_d\theta_d)^\top$。
GD 更新 $\theta_{t+1} = \theta_t - \eta\grad f(\theta_t)$ 在坐标 $i$ 上化为：
\begin{equation}\label{s2:coord}
  \theta_{t+1,i} = \theta_{t,i} - \eta \lambda_i \theta_{t,i}
  = (1 - \eta\lambda_i) \theta_{t,i}.
\end{equation}

递推 $T$ 步：
\begin{equation}\label{s2:evol}
  \boxed{\theta_{T,i} = (1 - \eta\lambda_i)^T \theta_{0,i}.}
\end{equation}

由于最优解为 $\theta^\star = 0$（因为 $f(\theta) = \frac12\theta^\top A\theta \ge 0$，
且 $A \succ 0$），误差由 $|1 - \eta\lambda_i|^T$ 控制。

\medskip
\noindent\textbf{第二步：最坏坐标的收敛因子。}

总体收敛速度由衰减最慢的坐标决定。定义
\begin{equation}
  \rho(\eta) = \max_{i=1,\ldots,d} |1 - \eta \lambda_i|.
\end{equation}

这是迭代矩阵 $I - \eta A$ 的谱半径。我们要找到 $\eta$ 最小化 $\rho(\eta)$。

令 $\lambda_{\min} = \min_i \lambda_i$，$\lambda_{\max} = \max_i \lambda_i$，
$\kappa = \lambda_{\max} / \lambda_{\min}$。函数 $|1 - \eta\lambda|$ 随 $\lambda$ 的变化如下图：

\begin{center}
\begin{itemize}
  \item 当 $\eta \lambda < 1$ 时，$1 - \eta\lambda > 0$。
  \item 当 $\eta \lambda > 1$ 时，$1 - \eta\lambda < 0$。
\end{itemize}
\end{center}

$\rho(\eta) = \max(|1-\eta\lambda_{\min}|, |1-\eta\lambda_{\max}|)$。

\medskip
\noindent\textbf{第三步：求最优 $\eta$。}

考虑以下三种情况：

\begin{itemize}
  \item 若 $\eta$ 很小（$\eta\lambda_{\max} < 1$），则两项皆正，
    $\rho(\eta) = 1 - \eta\lambda_{\min}$（最小值方向的衰减决定最坏情况）。
  \item 若 $\eta$ 较大（$\eta\lambda_{\min} > 1$），则两项皆负，
    $\rho(\eta) = \eta\lambda_{\max} - 1$（最大值方向决定最坏情况）。
  \item 最优 $\eta$ 在两者交叉处达到：
  \begin{equation}
    1 - \eta\lambda_{\min} = \eta\lambda_{\max} - 1.
  \end{equation}
\end{itemize}

解出：
\begin{equation}
  2 = \eta(\lambda_{\min} + \lambda_{\max})
  \implies \boxed{\eta^\star = \frac{2}{\lambda_{\min} + \lambda_{\max}}.}
\end{equation}

\medskip
\noindent\textbf{第四步：最优收敛因子。}

代入最优步长 $\eta^\star$：
\begin{equation}
  \begin{aligned}
    \rho(\eta^\star)
    &= 1 - \eta^\star \lambda_{\min}
     = 1 - \frac{2\lambda_{\min}}{\lambda_{\min} + \lambda_{\max}} \\[4pt]
    &= \frac{\lambda_{\min} + \lambda_{\max} - 2\lambda_{\min}}{\lambda_{\min} + \lambda_{\max}} \\[4pt]
    &= \frac{\lambda_{\max} - \lambda_{\min}}{\lambda_{\max} + \lambda_{\min}} \\[4pt]
    &= \boxed{\frac{\kappa - 1}{\kappa + 1}.}
  \end{aligned}
\end{equation}

因此，经过 $T$ 步 GD 后：
\begin{equation}
  \boxed{\norm{\theta_T}^2 \le \left(\frac{\kappa - 1}{\kappa + 1}\right)^{2T} \norm{\theta_0}^2.}
\end{equation}

对应的函数值收敛：
\begin{equation}
  f(\theta_T) \le \frac{\lambda_{\max}}{2} \left(\frac{\kappa - 1}{\kappa + 1}\right)^{2T} \norm{\theta_0}^2.
\end{equation}

\textbf{注记：} 当 $\eta = 1/\lambda_{\max}$ 时（较保守的步长），收敛因子为 $1 - 1/\kappa$。
两种步长的对比：
\[
  \frac{\kappa - 1}{\kappa + 1} = 1 - \frac{2}{\kappa + 1}
  \quad\text{vs}\quad
  1 - \frac{1}{\kappa}.
\]
当 $\kappa \gg 1$ 时，$\frac{\kappa-1}{\kappa+1} \approx 1 - \frac{2}{\kappa}$，
而 $1 - \frac{1}{\kappa} \approx 1 - \frac{1}{\kappa}$，
前者后者的两倍快（即其 $\log$ 衰减率约两倍）。
\end{proof*}

% ---------- 练习 3 ----------
\subsection{练习 3：条件数与梯度的几何}

\textbf{题目：} 研究 GD 在 $f(\theta) = \frac12\theta^\top Q\diag(1,\kappa)Q^\top\theta$ 上的轨迹。
分别考虑 $Q=I$ 和 $Q$ 为任意旋转的情况。

\begin{proof*}[解答]

\noindent\textbf{情况一：$Q=I$（坐标对齐，$A = \diag(1, \kappa)$）}

设 $\kappa = 10$。梯度为 $\grad f(\theta) = (\theta_1,\; \kappa\theta_2)^\top$。
GD 更新：
\begin{equation}
  \begin{pmatrix}\theta_1^{(t+1)} \\ \theta_2^{(t+1)}\end{pmatrix}
  = \begin{pmatrix}1 - \eta & 0 \\ 0 & 1 - \kappa\eta\end{pmatrix}
  \begin{pmatrix}\theta_1^{(t)} \\ \theta_2^{(t)}\end{pmatrix}.
\end{equation}

取 $\eta = 2/(1+\kappa) = 2/11 \approx 0.182$。则：
\begin{align}
  \theta_1^{(t)} &= \left(1 - \frac{2}{11}\right)^t \theta_1^{(0)}
  = \left(\frac{9}{11}\right)^t \theta_1^{(0)}, \\[4pt]
  \theta_2^{(t)} &= \left(1 - \frac{20}{11}\right)^t \theta_2^{(0)}
  = \left(-\frac{9}{11}\right)^t \theta_2^{(0)}.
\end{align}

$\theta_1$ 缓慢衰减（方向 1 是「慢方向」，曲率小），$\theta_2$ 快速衰减且符号翻转（方向 2 是「快方向」，曲率大）。
轨迹在 $(\theta_1,\theta_2)$ 平面中呈现「之字形」锯齿形下降，沿着慢方向（$\theta_1$ 轴）缓慢爬向原点。

\medskip
\noindent\textbf{情况二：$Q \neq I$（任意旋转）}

设 $A = Q\diag(1,\kappa)Q^\top$，其中 $Q$ 是正交矩阵（$QQ^\top = I$）。
令 $y = Q^\top\theta$（坐标变换到 Hessian 的特征基）。则：
\begin{equation}
  f(\theta) = \frac12\theta^\top A\theta
  = \frac12\theta^\top Q\diag(1,\kappa)Q^\top\theta
  = \frac12 y^\top \diag(1,\kappa) y = \frac12(y_1^2 + \kappa y_2^2).
\end{equation}

GD 在 $\theta$ 空间的更新：
\begin{equation}
  \theta_{t+1} = \theta_t - \eta A\theta_t
  = \theta_t - \eta Q\diag(1,\kappa)Q^\top\theta_t.
\end{equation}

在 $y$ 空间中：
\begin{equation}
  y_{t+1} = Q^\top\theta_{t+1}
  = Q^\top(\theta_t - \eta Q\diag(1,\kappa)Q^\top\theta_t)
  = y_t - \eta \diag(1,\kappa) y_t
  = (I - \eta\diag(1,\kappa)) y_t.
\end{equation}

因此\textbf{$y$ 空间中的递推与情况一完全一致}！GD 的动力学在旋转下完全等价。

\medskip
\noindent\textbf{为什么 GD 是旋转不变的？}

因为无约束 GD 的预条件器是 $I$（单位矩阵），且 $\norm{\cdot}_2$ 在正交变换下不变。
任意旋转只是重新标记了坐标轴，GD 的「下降方向 = 负梯度方向」这一几何事实不改变。

\medskip
\noindent\textbf{为什么对角预条件方法（如 AdaGrad / Adam）在旋转下表现不同？}

对角预条件器 $\diag(h_1, \dots, h_d)$ 只缩放坐标轴，不做旋转变换。
若 $A = Q\diag(1,\kappa)Q^\top$ 且 $Q$ 不是置换矩阵，则对角预条件器 $\diag(\sqrt{v_t})$并\textbf{不能}与 Hessian 的特征方向对齐。这意味着：

\begin{itemize}
  \item 在 $Q=I$ 时，对角预条件器可以完美适应 Hessian（只需对坐标 2 使用较小的学习率）。
  \item 在 $Q \neq I$ 时，对角预条件器无法去除坐标间的耦合——它的效果等价于在原始空间
    中使用一个对角的正定矩阵 $D$ 代替 $I$ 作为预条件器，即 $\theta_{t+1} \approx \theta_t - \eta D^{-1}\grad f(\theta_t)$。
    这改变了优化几何（从欧氏变为 Mahalanobis），但可能不足以应对 Hessian 的非对角结构。
\end{itemize}

这是为什么 Adam 等自适应方法在全连接层（特征维度间无自然排序）上的表现有时不如 SGD，而在嵌入/稀疏特征（特征维度有自然意义且 coordinate-aligned 的对角 Hessian 近似成立）上表现极佳。

\end{proof*}

% ---------- 练习 4 ----------
\subsection{练习 4：SGD 方差的实验测量}

\textbf{题目：} 在二维二次函数上向梯度添加高斯噪声模拟 SGD，研究最终误差与步长的关系。

\begin{proof*}[理论与实验分析]

\noindent\textbf{理论分析：}

设 $f(\theta) = \frac12\theta^\top A\theta$，$A \succ 0$。随机梯度 $g_t = A\theta_t + \xi_t$，
其中 $\xi_t \sim \mathcal{N}(0, \sigma^2 I)$。

由 SGD 在凸光滑条件下的分析（见 week1 讲义定理 5：SGD 在凸光滑条件下的收敛率）：
\begin{equation}
  \E[f(\bar\theta_T)] - f(\theta^\star) \le
  \frac{\norm{\theta_0}^2}{2\eta T} + \frac{\eta\sigma^2}{2}.
\end{equation}

当 $T \to \infty$ 时，第一项 $\to 0$，但第二项 $\eta\sigma^2/2$ 不随 $T$ 减小而消失。
这就是 SGD 的\textbf{稳态震荡}：
\begin{equation}
  \boxed{\lim_{T\to\infty} \E[f(\theta_T)] - f(\theta^\star) \approx \frac{\eta\sigma^2}{2}.}
\end{equation}

\noindent\textbf{更精确的分析（针对二次函数）：}

考虑一维情形 $f(\theta) = \frac{\lambda}{2}\theta^2$。SGD 更新：
\begin{equation}
  \theta_{t+1} = \theta_t - \eta(\lambda\theta_t + \xi_t) = (1 - \eta\lambda)\theta_t - \eta\xi_t.
\end{equation}

$\theta_t$ 是一个一阶自回归过程。当 $|1 - \eta\lambda| < 1$（即 $\eta < 2/\lambda$）时，
$\theta_t$ 趋于平稳分布，其平稳方差为：
\begin{equation}
  \Var(\theta_\infty) = \frac{\eta^2\sigma^2}{1 - (1 - \eta\lambda)^2}
  = \frac{\eta\sigma^2}{\lambda(2 - \eta\lambda)}.
\end{equation}

因此：
\begin{equation}
  \E[f(\theta_\infty)] = \frac{\lambda}{2}\E[\theta_\infty^2]
  = \frac{\lambda}{2} \cdot \frac{\eta\sigma^2}{\lambda(2 - \eta\lambda)}
  = \frac{\eta\sigma^2}{2(2 - \eta\lambda)}.
\end{equation}

取小步长 $\eta\lambda \ll 1$：
\begin{equation}
  \E[f(\theta_\infty)] \approx \frac{\eta\sigma^2}{4}.
\end{equation}

\noindent\textbf{最优步长的 trade-off：}

总误差（有限 $T$ 下）由两项组成：
\begin{equation}
  E(\eta) = \underbrace{\frac{\norm{\theta_0}^2}{2\eta T}}_{\text{收敛项（随 $\eta$ 增大而减小）}}
  + \underbrace{\frac{\eta\sigma^2}{2}}_{\text{稳态误差（随 $\eta$ 增大而增大）}}.
\end{equation}

令 $E'(\eta) = 0$：
\begin{equation}
  -\frac{\norm{\theta_0}^2}{2\eta^2 T} + \frac{\sigma^2}{2} = 0
  \implies \boxed{\eta_{\text{opt}} = \frac{\norm{\theta_0}}{\sigma\sqrt{T}}.}
\end{equation}

此时最优误差：
\begin{equation}
  \boxed{E(\eta_{\text{opt}}) = \frac{\norm{\theta_0}\sigma}{\sqrt{T}}.}
\end{equation}

这精确匹配了定理中的 $\mathcal{O}(1/\sqrt{T})$ 速率。

\medskip
\noindent\textbf{实验预期：}

\begin{itemize}
  \item 固定 $T = 1000$，变化 $\eta \in [10^{-3}, 1]$。
  \item $\eta$ 很小时：收敛项大，误差大。
  \item $\eta$ 在最优值附近：误差最小。
  \item $\eta$ 较大时：稳态噪声主导，误差大。
  \item $\eta \ge 2/\lambda_{\max}$ 时：发散。
\end{itemize}

典型的「最终误差 vs $\eta$」曲线呈 U 形，底部对应最优步长。

\end{proof*}

% ---------- 练习 5 ----------
\subsection{练习 5：Gradient Descent + Momentum 的特征分析}

\textbf{题目：} 在一维二次函数 $f(\theta) = \frac{\lambda}{2}\theta^2$ 上，
分析 GD+Momentum 的状态空间方程，找出使谱半径最小化的 $(\eta, \beta)$。

\begin{proof*}[解答]

\noindent\textbf{第一步：建立状态空间方程。}

一维二次函数：$\grad f(\theta) = \lambda\theta$。

Momentum 更新（Polyak 重球法）：
\begin{equation}\label{s5:momentum}
  \begin{aligned}
    d_t &= \beta d_{t-1} + \lambda\theta_t, \\
    \theta_{t+1} &= \theta_t - \eta d_t.
  \end{aligned}
\end{equation}

将 $d_t$ 代入 $\theta_{t+1}$ 的表达式：
\begin{equation}
  \theta_{t+1} = \theta_t - \eta(\beta d_{t-1} + \lambda\theta_t) = (1 - \eta\lambda)\theta_t - \eta\beta d_{t-1}.
\end{equation}

定义状态向量 $x_t = (\theta_t,\; d_{t-1})^\top$。则：
\begin{equation}
  \begin{aligned}
    \theta_{t+1} &= (1 - \eta\lambda)\theta_t - \eta\beta d_{t-1}, \\
    d_t &= \beta d_{t-1} + \lambda\theta_t.
  \end{aligned}
\end{equation}

写为矩阵形式：
\begin{equation}
  \boxed{
  \begin{pmatrix} \theta_{t+1} \\ d_t \end{pmatrix}
  = \underbrace{\begin{pmatrix} 1 - \eta\lambda & -\eta\beta \\ \lambda & \beta \end{pmatrix}}_{M}
  \begin{pmatrix} \theta_t \\ d_{t-1} \end{pmatrix}.
  }
\end{equation}

\noindent\textbf{第二步：求特征值。}

矩阵 $M$ 的特征多项式：
\begin{equation}
  \det(M - \mu I) =
  \begin{vmatrix}
    1 - \eta\lambda - \mu & -\eta\beta \\
    \lambda & \beta - \mu
  \end{vmatrix} = 0.
\end{equation}

展开：
\begin{equation}
  (1 - \eta\lambda - \mu)(\beta - \mu) - (-\eta\beta)(\lambda) = 0.
\end{equation}
\begin{equation}
  (1 - \eta\lambda - \mu)(\beta - \mu) + \eta\beta\lambda = 0.
\end{equation}
\begin{equation}
  \mu^2 - (1 - \eta\lambda + \beta)\mu + \beta(1 - \eta\lambda) + \eta\beta\lambda = 0.
\end{equation}
\begin{equation}
  \mu^2 - (1 - \eta\lambda + \beta)\mu + \beta = 0.
\end{equation}

两个特征值为：
\begin{equation}
  \boxed{\mu_{1,2} = \frac{(1 - \eta\lambda + \beta) \pm \sqrt{(1 - \eta\lambda + \beta)^2 - 4\beta}}{2}.}
\end{equation}

\noindent\textbf{第三步：最小化谱半径。}

谱半径 $\rho(M) = \max(|\mu_1|, |\mu_2|)$。我们希望最小化 $\rho(M)$。

有两种情况：

\begin{itemize}
  \item \textbf{实特征值（判别式 $\ge 0$）}：$\mu_{1,2}$ 皆为正实数（假设 $\beta \ge 0$, $\eta$ 适当），
    则 $\rho(M) = \mu_{\max}$。
  \item \textbf{复特征值（判别式 $< 0$）}：$\mu_{1,2}$ 为共轭复数，则 $|\mu_1| = |\mu_2| = \sqrt{\beta}$。
\end{itemize}

在复特征值情形下，谱半径为 $\sqrt{\beta}$，与 $\eta$ 无关（只要判别式 $<0$ 的条件成立）。
这意味著只要步长选在合适的区间内，收敛速度只由动量系数 $\beta$ 决定。

判别式 $< 0$ 的条件：
\begin{equation}
  (1 - \eta\lambda + \beta)^2 < 4\beta.
\end{equation}
即：
\begin{equation}
  -2\sqrt{\beta} < 1 - \eta\lambda + \beta < 2\sqrt{\beta}.
\end{equation}
整理关于 $\eta\lambda$ 的条件：
\begin{equation}
  (1 - \sqrt{\beta})^2 < \eta\lambda < (1 + \sqrt{\beta})^2.
\end{equation}

在此区间内，$\rho(M) = \sqrt{\beta}$。我们可以选择 $\eta$ 使判别式为零（临界阻尼临界点，此时收敛最快）：
\begin{equation}
  (1 - \eta\lambda + \beta)^2 = 4\beta \implies \eta\lambda = (1 - \sqrt{\beta})^2.
\end{equation}
或 $\eta\lambda = (1 + \sqrt{\beta})^2$（后者对应的 $\eta$ 更大，通常不选）。

取 $\eta\lambda = (1 - \sqrt{\beta})^2$，此实为「临界复特征值」的边界。

\medskip
\noindent\textbf{第四步：推广到多维和加速因子。}

对于多维二次函数 $f(\theta) = \frac12\theta^\top A\theta$，在 $A$ 的特征基下解耦为
$d$ 个独立的过程。令 $\kappa = \lambda_{\max}/\lambda_{\min}$。最优参数选择为：
\begin{equation}
  \eta^\star = \frac{4}{(\sqrt{\lambda_{\max}} + \sqrt{\lambda_{\min}})^2},
  \qquad
  \beta^\star = \left(\frac{\sqrt{\kappa} - 1}{\sqrt{\kappa} + 1}\right)^2.
\end{equation}

在此参数下，收敛因子（谱半径）为：
\begin{equation}
  \boxed{\rho^\star = \sqrt{\beta^\star} = \frac{\sqrt{\kappa} - 1}{\sqrt{\kappa} + 1}.}
\end{equation}

\medskip
\noindent\textbf{对比：GD vs Momentum 的收敛因子}

\begin{center}
\begin{tabular}{lcc}
\toprule
方法 & 收敛因子（每步） & $\kappa=100$ 时的数值 \\
\midrule
GD（$\eta = 1/\lambda_{\max}$） & $1 - \frac{1}{\kappa}$ & 0.99 \\
GD（$\eta = 2/(\lambda_{\min} + \lambda_{\max})$） & $\frac{\kappa-1}{\kappa+1}$ & 0.9802 \\
Momentum（最优参数） & $\frac{\sqrt{\kappa}-1}{\sqrt{\kappa}+1}$ & 0.8182 \\
\bottomrule
\end{tabular}
\end{center}

\medskip
\textbf{加速因子的定量解释：}

当 $\kappa \gg 1$ 时：
\begin{align}
  \rho_{\text{GD}} &= \frac{\kappa-1}{\kappa+1}
    \approx 1 - \frac{2}{\kappa}, \\[4pt]
  \rho_{\text{Momentum}} &= \frac{\sqrt{\kappa}-1}{\sqrt{\kappa}+1}
    \approx 1 - \frac{2}{\sqrt{\kappa}}.
\end{align}

达到误差 $\eps$ 所需的迭代数：
\begin{align}
  T_{\text{GD}} &\approx \frac{\kappa}{2} \log\frac{1}{\eps}, \\[4pt]
  T_{\text{Momentum}} &\approx \frac{\sqrt{\kappa}}{2} \log\frac{1}{\eps}.
\end{align}

当 $\kappa = 10000$ 时，Momentum 将迭代数减少约 $100\times$。这就是「加速」的含义。

\end{proof*}

% ======================================================================
\section{工程练习}
% ======================================================================

% ---------- 工程任务 1 ----------
\subsection{工程任务 1：实现并可视化 GD}

\noindent\textbf{目标：} 在 $\kappa=100$ 的二次函数上验证 GD 的收敛行为。

\begin{lstlisting}[language=Python, caption={GD 在不同步长下的收敛行为}]
import numpy as np
import matplotlib.pyplot as plt

# ---- 问题定义 ----
d = 2
A = np.diag([1.0, 100.0])           # kappa = 100
theta0 = np.array([10.0, 1.0])      # 初始点
theta_star = np.zeros(d)            # 最优解

def f(theta):
    return 0.5 * theta @ A @ theta

def grad_f(theta):
    return A @ theta

# ---- GD 运行 ----
L = np.max(np.diag(A))               # lambda_max = 100
etas = [0.0001, 0.001, 0.01, 0.019, 0.02, 0.021]
T = 200

fig, (ax1, ax2) = plt.subplots(1, 2, figsize=(14, 5))

for eta in etas:
    theta = theta0.copy()
    losses = []
    thetas = [theta.copy()]

    for t in range(T):
        g = grad_f(theta)
        theta = theta - eta * g
        losses.append(f(theta))
        thetas.append(theta.copy())

    # 左图: loss curve
    factor = eta * L
    label = f"$\\eta={eta}$, $\\eta L={factor:.2f}$"
    ax1.semilogy(losses, label=label, alpha=0.8)
    ax1.axhline(1e-15, color='gray', ls='--', alpha=0.3)

    # 右图: trajectory (只画不发散的)
    thetas = np.array(thetas)
    if losses[-1] < 1e5:  # 未发散
        ax2.plot(thetas[:, 0], thetas[:, 1], '.-',
                 alpha=0.5, markersize=2, label=label)

ax1.set_xlabel('迭代次数')
ax1.set_ylabel('$f(\\theta_t)$ (log scale)')
ax1.set_title('GD 在不同步长下的收敛 ($\\kappa=100$)')
ax1.legend(fontsize=7, loc='lower left')
ax1.grid(True, alpha=0.3)

ax2.set_xlabel('$\\theta_1$')
ax2.set_ylabel('$\\theta_2$')
ax2.set_title('GD 轨迹 ($\\kappa=100$)')
ax2.legend(fontsize=7)
ax2.grid(True, alpha=0.3)
ax2.set_xlim(-12, 12)
ax2.set_ylim(-1.5, 1.5)

plt.tight_layout()
plt.savefig('task1_gd_convergence.pdf', dpi=150)
plt.show()
\end{lstlisting}

\textbf{关键观察：}

\begin{itemize}
  \item $\eta L < 1$ 时单调下降；$\eta L = 1$ 时最快下降；$\eta L > 2$ 时发散。
  \item 轨迹图中，GD 在慢方向（$\theta_1$ 轴）进展缓慢，在快方向（$\theta_2$ 轴）迅速收敛。
  \item 由于 $\kappa=100$，$\theta_1$ 的衰减因子约为 $1 - \eta$，而 $\theta_2$ 的衰减因子约为 $1 - 100\eta$。
    在 $\eta = 1/L = 0.01$ 时，前者为 0.99（极慢），后者为 0（一步收敛）——这正是「之字形」轨迹的数学原因。
  \item 最优步长 $\eta = 2/(1+100) \approx 0.0198$ 时，两个方向的衰减因子绝对值相等（$99/101 \approx 0.98$），
    但代价是 $\theta_2$ 方向会产生符号翻转（overshoot）。
\end{itemize}

% ---------- 工程任务 2 ----------
\subsection{工程任务 2：SGD 的噪声与方差}

\noindent\textbf{目标：} 在二次函数上验证 SGD 的 $\mathcal{O}(1/\sqrt{T})$ 收敛率。

\begin{lstlisting}[language=Python, caption={SGD 噪声方差与收敛率验证}]
import numpy as np
import matplotlib.pyplot as plt

# ---- 问题定义 ----
A = np.diag([1.0, 100.0])
theta0 = np.array([10.0, 1.0])

def f(theta):
    return 0.5 * theta @ A @ theta

def grad_f(theta):
    return A @ theta

# ---- (A) 不同 sigma 下的 SGD ----
eta = 1.0 / 100.0   # 固定步长 eta = 1/L
sigmas = [0.1, 1.0, 3.0]
T = 500
n_runs = 10
rng = np.random.RandomState(42)

fig, axes = plt.subplots(1, 2, figsize=(14, 5))
ax1, ax2 = axes

colors = ['blue', 'orange', 'red']

for si, sigma in enumerate(sigmas):
    all_losses = np.zeros((n_runs, T))

    for run in range(n_runs):
        theta = theta0.copy()
        for t in range(T):
            noise = sigma * rng.randn(2)
            g = grad_f(theta) + noise
            theta = theta - eta * g
            all_losses[run, t] = f(theta)

    # 均值与标准差带
    mean_loss = all_losses.mean(axis=0)
    std_loss = all_losses.std(axis=0)
    ax1.semilogy(mean_loss, color=colors[si],
                 label=f'$\\sigma={sigma}$')
    ax1.fill_between(range(T),
                     np.maximum(mean_loss - std_loss, 1e-15),
                     mean_loss + std_loss,
                     alpha=0.2, color=colors[si])

# 加 GD 对照
theta = theta0.copy()
gd_losses = []
for t in range(T):
    g = grad_f(theta)
    theta = theta - eta * g
    gd_losses.append(f(theta))
ax1.semilogy(gd_losses, 'k--', label='GD (无噪声)', alpha=0.8)

ax1.set_xlabel('迭代次数 $t$')
ax1.set_ylabel('$f(\\theta_t)$ (log scale)')
ax1.set_title('SGD 在不同噪声水平下的收敛')
ax1.legend(fontsize=9)
ax1.grid(True, alpha=0.3)

# ---- (B) 收敛率验证: log(error) vs log(T) ----
T_large = 5000
eta = 0.01
sigma = 1.0

theta = theta0.copy()
errors = []
for t in range(T_large):
    noise = sigma * rng.randn(2)
    g = grad_f(theta) + noise
    theta = theta - eta * g
    if (t + 1) % 50 == 0:
        errors.append((t + 1, f(theta)))

errors = np.array(errors)
ax2.loglog(errors[:, 0], errors[:, 1], 'b-', alpha=0.8,
           label='SGD ($\\eta=0.01$, $\\sigma=1$)')

# 参考线: O(1/sqrt(T)) 和 O(1/T)
t_vals = errors[:, 0]
ref_sqrt = 10 / np.sqrt(t_vals)       # O(1/sqrt(T))
ref_linear = 100 / t_vals              # O(1/T)

ax2.loglog(t_vals, ref_sqrt, 'r--', alpha=0.7, label='$\\mathcal{O}(1/\\sqrt{T})$')
ax2.loglog(t_vals, ref_linear, 'g--', alpha=0.7, label='$\\mathcal{O}(1/T)$')

ax2.set_xlabel('迭代次数 $T$ (log scale)')
ax2.set_ylabel('$f(\\theta_T)$ (log scale)')
ax2.set_title('SGD 收敛率验证')
ax2.legend(fontsize=9)
ax2.grid(True, alpha=0.3)

plt.tight_layout()
plt.savefig('task2_sgd_noise.pdf', dpi=150)
plt.show()

# ---- (C) 稳态误差 vs eta ----
sigmas_test = [0.5, 1.0, 2.0]
etas_test = np.logspace(-3, -0.3, 30)
T_warm = 5000

fig3, ax3 = plt.subplots(figsize=(8, 5))

for sigma in sigmas_test:
    final_errors = []
    for eta in etas_test:
        # 检查稳定性
        if eta >= 2.0 / 100.0:
            final_errors.append(np.nan)
            continue
        theta = np.zeros(2)
        # 预热
        for t in range(T_warm):
            noise = sigma * rng.randn(2)
            g = grad_f(theta) + noise
            theta = theta - eta * g
        # 测量稳态
        steady = 0.0
        for t in range(500):
            noise = sigma * rng.randn(2)
            g = grad_f(theta) + noise
            theta = theta - eta * g
            steady += f(theta)
        final_errors.append(steady / 500)

    final_errors = np.array(final_errors)
    valid = ~np.isnan(final_errors)
    ax3.loglog(etas_test[valid], final_errors[valid], 'o-',
               label=f'$\\sigma={sigma}$', alpha=0.8)

    # 理论线: eta * sigma^2 / 4
    theory = etas_test * sigma**2 / 4
    ax3.loglog(etas_test, theory, '--', alpha=0.4,
               label=f'理论: $\\eta\\sigma^2/4$' if sigma == sigmas_test[0] else '')

ax3.set_xlabel('步长 $\\eta$ (log scale)')
ax3.set_ylabel('稳态误差 (log scale)')
ax3.set_title('SGD 稳态误差 vs 步长 (验证 $\\propto \\eta\\sigma^2$)')
ax3.legend(fontsize=9)
ax3.grid(True, alpha=0.3)

plt.tight_layout()
plt.savefig('task2_steady_state.pdf', dpi=150)
plt.show()

print("=== 关键发现 ===")
print("1. SGD 的误差在 log-log 图上的斜率约为 -0.5，验证了 O(1/sqrt(T)) 速率。")
print("2. 噪声越大 (sigma 大)，收敛越慢且稳态误差越高。")
print("3. 稳态误差正比于 eta * sigma^2，与理论一致。")
\end{lstlisting}

\textbf{关键观察：}

\begin{itemize}
  \item 左图展示了不同 $\sigma$ 下 SGD 的均值和标准差带。$\sigma$ 越大，收敛后的震荡幅度越大。
  \item 右图（log-log）中 SGD 的斜率约为 $-1/2$，而 GD 的参考线斜率为 $-1$——直观验证了 SGD 的慢 $\sqrt{T}$ 因子。
  \item 最终误差 vs $\eta$ 的曲线验证了理论：在小 $\eta$ 时受收敛速度限制，在大 $\eta$ 时受稳态噪声限制，
    存在一个最优 $\eta$ 平衡两者。
\end{itemize}

% ---------- 工程任务 3 ----------
\subsection{工程任务 3：Momentum 与 Nesterov 的加速效果}

\noindent\textbf{目标：} 在 $\kappa=1000$ 的二次函数上比较 GD、Momentum、Nesterov。

\begin{lstlisting}[language=Python, caption={GD vs Momentum vs Nesterov 加速对比}]
import numpy as np
import matplotlib.pyplot as plt

# ---- 问题定义 ----
kappa = 1000
A = np.diag([1.0, kappa])              # kappa = 1000
theta0 = np.array([10.0, 1.0])
theta_star = np.zeros(2)

def f(theta):
    return 0.5 * theta @ A @ theta

def grad_f(theta):
    return A @ theta

lambda_min = 1.0
lambda_max = kappa
L = lambda_max
mu = lambda_min

# ---- 最优参数 ----
# GD: eta = 2/(lambda_min + lambda_max)
eta_gd = 2.0 / (lambda_min + lambda_max)

# Momentum (Polyak): 最优参数
beta_momentum = ((np.sqrt(kappa) - 1) / (np.sqrt(kappa) + 1))**2
eta_momentum = (2 - 2 * beta_momentum) / lambda_max

# Nesterov: 标准参数
eta_nest = 1.0 / L
beta_nest = (np.sqrt(kappa) - 1) / (np.sqrt(kappa) + 1)

T = 2000

# ---- 运行三种优化器 ----
def run_gd(theta0, eta, T):
    theta = theta0.copy()
    losses, thetas = [], [theta.copy()]
    for t in range(T):
        theta = theta - eta * grad_f(theta)
        losses.append(f(theta))
        thetas.append(theta.copy())
    return np.array(losses), np.array(thetas)

def run_momentum(theta0, eta, beta, T):
    theta = theta0.copy()
    d = np.zeros_like(theta)
    losses, thetas = [], [theta.copy()]
    for t in range(T):
        g = grad_f(theta)
        d = beta * d + g
        theta = theta - eta * d
        losses.append(f(theta))
        thetas.append(theta.copy())
    return np.array(losses), np.array(thetas)

def run_nesterov(theta0, eta, beta, T):
    theta = theta0.copy()
    d = np.zeros_like(theta)
    losses, thetas = [], [theta.copy()]
    for t in range(T):
        # Nesterov: 在前瞻点计算梯度
        lookahead = theta - eta * beta * d
        g = grad_f(lookahead)
        d = beta * d + g
        theta = theta - eta * d
        losses.append(f(theta))
        thetas.append(theta.copy())
    return np.array(losses), np.array(thetas)

loss_gd, traj_gd = run_gd(theta0, eta_gd, T)
loss_mom, traj_mom = run_momentum(theta0, eta_momentum, beta_momentum, T)
loss_nest, traj_nest = run_nesterov(theta0, eta_nest, beta_nest, T)

# ---- 绘图 ----
fig, axes = plt.subplots(1, 2, figsize=(14, 5))

# 左图: 收敛曲线
ax1 = axes[0]
ax1.semilogy(loss_gd, 'b-', alpha=0.8, label=f'GD ($\\eta=2/(\\mu+L)$)')
ax1.semilogy(loss_mom, 'orange', alpha=0.8,
             label=f'Momentum ($\\beta={beta_momentum:.3f}$)')
ax1.semilogy(loss_nest, 'g-', alpha=0.8,
             label=f'Nesterov ($\\beta={beta_nest:.3f}$)')

# 理论衰减参考线
t = np.arange(T)
rho_gd = (kappa - 1) / (kappa + 1)
rho_acc = (np.sqrt(kappa) - 1) / (np.sqrt(kappa) + 1)
ref_gd = f(theta0) * rho_gd**(2*t)
ref_acc = f(theta0) * rho_acc**t
ax1.semilogy(ref_gd, 'b--', alpha=0.3, label=f'理论: GD $O((\\kappa-1)/(\\kappa+1))^{{2T}}$')
ax1.semilogy(ref_acc, 'r--', alpha=0.3, label=f'理论: Accel $O((\\sqrt{{\\kappa}}-1)/(\\sqrt{{\\kappa}}+1))^T$')

ax1.set_xlabel('迭代次数')
ax1.set_ylabel('$f(\\theta_t)$ (log scale)')
ax1.set_title(f'GD vs Momentum vs Nesterov ($\\kappa={kappa}$)')
ax1.legend(fontsize=8)
ax1.grid(True, alpha=0.3)

# 右图: 轨迹
ax2 = axes[1]
labels = [('GD', traj_gd, 'b'), ('Momentum', traj_mom, 'orange'),
          ('Nesterov', traj_nest, 'g')]
for name, traj, color in labels:
    ax2.plot(traj[:80, 0], traj[:80, 1], '-', color=color,
             alpha=0.7, linewidth=1, label=name)
    ax2.plot(traj[0, 0], traj[0, 1], 'o', color=color, markersize=6)

ax2.plot(0, 0, 'r*', markersize=12, label='Optimum')
ax2.set_xlabel('$\\theta_1$')
ax2.set_ylabel('$\\theta_2$')
ax2.set_title(f'前 80 步轨迹对比 ($\\kappa={kappa}$)')
ax2.legend(fontsize=8)
ax2.grid(True, alpha=0.3)

plt.tight_layout()
plt.savefig('task3_acceleration.pdf', dpi=150)
plt.show()

# ---- 加速因子验证 ----
tol = 1e-10
for name, losses in [('GD', loss_gd), ('Momentum', loss_mom),
                      ('Nesterov', loss_nest)]:
    steps = np.where(losses < tol)[0]
    if len(steps) > 0:
        print(f"{name}: 达到 {tol:.0e} 需要 {steps[0]} 步")
    else:
        print(f"{name}: 在 {T} 步内未达到 {tol:.0e}")

print(f"\n理论加速比: sqrt(kappa) = {np.sqrt(kappa):.1f}")
print(f"GD 理论收敛因子: {rho_gd:.6f}")
print(f"加速理论收敛因子: {rho_acc:.6f}")
\end{lstlisting}

\textbf{关键观察：}

\begin{itemize}
  \item GD 的衰减因子约为 $1 - 2/\kappa = 0.998$（极慢！），2000 步后误差仅下降约 $2\times$。
  \item Momentum 和 Nesterov 的衰减因子约为 $1 - 2/\sqrt{\kappa} \approx 0.939$，
    相同迭代数下的误差下降远快于 GD。
  \item 轨迹图（右）展示 Momentum/Nesterov 的「过冲」现象——沿快方向（$\theta_2$ 轴）越过最优点后折返，
    与 GD 沿慢方向的「之字形」形成鲜明对比。Momentum 利用了惯性在快方向上做了更大的步长并快速定位到谷底。
  \item Nesterov 方法在理论上与 Momentum 有相同的渐近速率，但「前瞻梯度」减少了过冲振幅，
    在实际中通常比标准 Momentum 稍快且更稳定。
\end{itemize}

% ---------- 工程任务 4 ----------
\subsection{工程任务 4：学习率调度对比}

\noindent\textbf{目标：} 在 logistic regression 上比较不同的学习率调度策略。

\begin{lstlisting}[language=Python, caption={学习率调度策略对比}]
import numpy as np
import matplotlib.pyplot as plt

# ---- 生成数据 ----
np.random.seed(42)
n, d = 2000, 50
X = np.random.randn(n, d)
theta_true = np.random.randn(d)
# logit: p(y=1|x) = sigmoid(x^T theta_true)
prob = 1.0 / (1.0 + np.exp(-X @ theta_true))
y = (np.random.rand(n) < prob).astype(np.float64)

def sigmoid(z):
    return 1.0 / (1.0 + np.exp(-np.clip(z, -50, 50)))

def logistic_loss(theta, X, y):
    p = sigmoid(X @ theta)
    eps = 1e-12
    return -np.mean(y * np.log(p + eps) + (1 - y) * np.log(1 - p + eps))

def logistic_grad(theta, X, y):
    return X.T @ (sigmoid(X @ theta) - y) / n

def sgd_with_schedule(theta0, X, y, T, lr_schedule):
    """运行 SGD 并返回 loss 历史。lr_schedule(t) 返回第 t 步的学习率。"""
    theta = theta0.copy()
    losses = []
    rng = np.random.RandomState(123)
    for t in range(T):
        # mini-batch
        idx = rng.choice(n, size=32, replace=False)
        g = logistic_grad(theta, X[idx], y[idx])
        eta = lr_schedule(t)
        theta = theta - eta * g
        if t % 20 == 0:
            losses.append((t, logistic_loss(theta, X, y)))
    return np.array(losses)

# ---- 调度定义 ----
T_total = 3000

# 1. 固定步长
def fixed(eta0=1.0):
    return lambda t: eta0

# 2. Warmup + Cosine Decay
def warmup_cosine(eta_max=1.0, eta_min=1e-3, Tw=300, T=T_total):
    def schedule(t):
        if t < Tw:
            return eta_max * t / Tw
        progress = (t - Tw) / (T - Tw)
        return eta_min + (eta_max - eta_min) * 0.5 * (1 + np.cos(np.pi * progress))
    return schedule

# 3. Step Decay
def step_decay(eta0=1.0, gamma=0.5, step_size=600):
    def schedule(t):
        return eta0 * (gamma ** (t // step_size))
    return schedule

# 4. 多项式衰减
def poly_decay(eta0=1.0, p=2.0, T=T_total):
    def schedule(t):
        return eta0 * (1 - t / T) ** p
    return schedule

# ---- 运行实验 ----
theta0 = np.zeros(d)
schedules = {
    'Fixed ($\\eta=1.0$)': fixed(1.0),
    'Fixed ($\\eta=0.1$)': fixed(0.1),
    'Warmup+Cosine': warmup_cosine(),
    'Step Decay ($\\gamma=0.5$)': step_decay(),
    'Poly Decay ($p=2$)': poly_decay(),
}

fig, (ax1, ax2) = plt.subplots(1, 2, figsize=(14, 5))
colors = plt.cm.tab10(np.linspace(0, 1, len(schedules)))

for (name, lr_fn), color in zip(schedules.items(), colors):
    result = sgd_with_schedule(theta0, X, y, T_total, lr_fn)
    ax1.semilogy(result[:, 0], result[:, 1], color=color,
                 label=name, alpha=0.8)

    # 画学习率曲线
    ts = np.arange(T_total)
    lrs = [lr_fn(t) for t in ts]
    ax2.plot(ts, lrs, color=color, label=name, alpha=0.8)

ax1.set_xlabel('迭代次数')
ax1.set_ylabel('Loss (log scale)')
ax1.set_title('不同学习率调度下的 SGD 收敛')
ax1.legend(fontsize=7)
ax1.grid(True, alpha=0.3)

ax2.set_xlabel('迭代次数')
ax2.set_ylabel('学习率 $\\eta_t$')
ax2.set_title('学习率调度曲线')
ax2.legend(fontsize=7)
ax2.grid(True, alpha=0.3)

plt.tight_layout()
plt.savefig('task4_lr_schedule.pdf', dpi=150)
plt.show()

# ---- 最终 loss 统计 ----
print("=== 最终 Loss (3000 步) ===")
for name, lr_fn in schedules.items():
    res = sgd_with_schedule(theta0, X, y, T_total, lr_fn)
    print(f"{name:30s}: {res[-1, 1]:.6f}")
\end{lstlisting}

\textbf{关键观察：}

\begin{itemize}
  \item \textbf{固定步长 $\eta=1.0$}：初期进展快，但后期在最优解附近剧烈震荡，最终 loss 较高。
  \item \textbf{固定步长 $\eta=0.1$}：初期进展慢，但最终收敛到较低水平——体现了收敛速度与稳态精度的 trade-off。
  \item \textbf{Warmup+Cosine}：综合了两者的优点——早期有足够的步长快速进展，后期逐步衰减实现精准收敛。
    通常在实践中表现最佳。
  \item \textbf{Step Decay}：在每个衰减点出现「阶梯式」的 loss 下降，体现为 loss 曲线上的不连续跳跃。
    这种方法简单但需要预先指定衰减时间点，对问题敏感。
  \item \textbf{多项式衰减（$p=2$）}：平滑地将学习率递减到 0，最终精度较高，但衰减速度可能过快导致
    训练早期结束。
\end{itemize}

\noindent\textbf{最佳实践建议}：在深度学习训练中，Warmup + Cosine Decay 是当前最广泛使用的组合。
Warmup 防止训练初期的不稳定，Cosine decay 提供了平滑且有效的学习率衰减。

% ======================================================================
\section{总结：第一周核心收获}
% ======================================================================

\begin{enumerate}[label=(\arabic*)]
  \item \textbf{下降引理}是所有 GD 收敛分析的起点——它利用 $L$-光滑性将函数值下降与梯度范数联系起来。
  \item \textbf{条件数 $\kappa = L/\mu$} 直接决定了 GD 在强凸函数上的收敛速度。
    GD 在 $\kappa$ 很大时极度缓慢，这是所有后续加速方法存在的根本原因。
  \item \textbf{SGD 的 $\mathcal{O}(1/\sqrt{T})$ 速率}源于随机梯度的方差 $\sigma^2$。
    方差越大，最优步长越小，收敛越慢。
  \item \textbf{Momentum/Nesterov 加速}将有效条件数从 $\kappa$ 降为 $\sqrt{\kappa}$，
    其本质是通过历史梯度的指数移动平均「放大」低频（持续下降的）方向的分量。
  \item \textbf{学习率调度}是 SGD 在实践中不可或缺的组件：Warmup 防止初期发散，
    衰减策略保证后期精确收敛。
\end{enumerate}

\end{document}
